\documentclass[border=2pt,tikz]{standalone}
% Shared TikZ/pgfplots style for every figure in the manuscript.
\usepackage[T1]{fontenc}
\usepackage{mathpazo}
\usepackage{amsmath,amssymb}
\usepackage{tikz}
\usepackage{pgfplots}
\usepackage{pgfplotstable}
\pgfplotsset{compat=1.18}
\usetikzlibrary{arrows.meta,positioning,calc,backgrounds,fit,shapes.geometric}
\usepgfplotslibrary{fillbetween}

% --- palette: colour-blind safe, legible in greyscale -------------------
\definecolor{cHyd}{HTML}{1B6CA8}   % hydride / primary blue
\definecolor{cPro}{HTML}{C1492E}   % proton  / primary red
\definecolor{cAcc}{HTML}{7B5AA6}   % accent  / purple
\definecolor{cGold}{HTML}{B8860B}  % accent  / gold
\definecolor{cGrey}{HTML}{6E6E6E}
\definecolor{cInk}{HTML}{1A1A1A}
\definecolor{cBandR}{HTML}{C1492E}
\definecolor{cBandG}{HTML}{4F7942}

% --- base axis ----------------------------------------------------------
\pgfplotsset{
  icig/.style={
    width=68mm, height=52mm,
    scale only axis,
    axis line style={cInk, line width=0.45pt},
    tick style={cInk, line width=0.4pt},
    tick align=outside, tick pos=left,
    label style={font=\small, color=cInk},
    tick label style={font=\footnotesize, color=cInk},
    title style={font=\small\bfseries, color=cInk, align=left},
    legend cell align=left,
    legend style={draw=none, fill=none, font=\scriptsize,
                  inner sep=1pt, row sep=0.5pt},
    grid=both,
    major grid style={cGrey!22, line width=0.28pt},
    minor grid style={cGrey!12, line width=0.2pt},
  },
  % marker shorthands
  ptA/.style={mark=*, mark size=1.7pt, thick,
              mark options={fill=cHyd, draw=white, line width=0.4pt}},
  ptB/.style={mark=*, mark size=1.7pt, thick,
              mark options={fill=cPro, draw=white, line width=0.4pt}},
  hollowB/.style={mark=o, mark size=1.7pt, thick,
              mark options={draw=cPro, line width=0.7pt}},
}

% panel tag, e.g. \panel{a}
\newcommand{\panel}[1]{\textbf{(\textsf{#1})}}

\usetikzlibrary{arrows.meta,positioning,calc}
\begin{document}
\begin{tikzpicture}
% Worked-example figure, full text width.  Panel B is computed by
% analysis/export_fig_example.py; panel A is a schematic.
% Generated by analysis/export_fig_example.py; do not edit.
\def\fxGamma{3.348866}
\def\fxTop{5.00}
\def\fxYadhE{3.583716}
\def\fxLightE{2.713877}

\pgfplotsset{
  ov/.style={icig, scale only axis,
    label style={font=\fontsize{8.0}{9.6}\selectfont, color=cInk},
    tick label style={font=\fontsize{7.2}{8.6}\selectfont, color=cInk},
    title style={font=\fontsize{8.6}{10.4}\selectfont\bfseries, color=cInk,
                 align=left}},
}
\tikzset{
  ann/.style={font=\fontsize{7.2}{8.6}\selectfont, color=cInk, align=left,
              inner sep=0pt},
  annc/.style={font=\fontsize{7.2}{8.6}\selectfont, color=cInk, align=center,
              inner sep=0pt},
  hd/.style={font=\fontsize{8.6}{10.4}\selectfont\bfseries, color=cInk,
             anchor=north west, inner sep=0pt},
  box/.style={draw=#1, fill=#1!7, line width=0.45pt, rounded corners=1.5pt,
              inner sep=3pt, font=\fontsize{7.2}{8.8}\selectfont, text=cInk,
              align=left},
  setarr/.style={line width=1.5pt, -{Latex[length=2.0mm,width=1.9mm]}},
}
\hyphenpenalty=10000 \exhyphenpenalty=10000

% =====================================================================
% (A) the experiment and what is unknown
% =====================================================================
\node[hd] at (0,0) {(\textit{\textbf{A}})\ \ the experiment};

\node[anchor=north west, inner sep=0pt, font=\fontsize{7.6}{9}\selectfont]
  (rx) at (0.05,-0.62)
  {$\mathrm{E}+\mathrm{S}_L\;
    \underset{k_{\mathrm{off}}}{\overset{k_{\mathrm{on}}}{\rightleftharpoons}}\;
    \mathrm{ES}_L\;\xrightarrow{\;k_L\;}\;\mathrm{E}+\mathrm{P}$};
\node[ann, anchor=north west, text width=50mm] (cmp)
  at ($(rx.south west)+(0,-0.16)$)
  {H competes with T in one assay and D with T in another. Both assays use the same tritiated substrate in this illustration.};

\node[box=cHyd, anchor=north west, text width=46mm] (meas)
  at ($(cmp.south west)+(0,-0.20)$)
  {\textbf{measured}\quad $\KHT$, $\KDT$\\[1pt]
   isotope effects on $V/K$ for H and for D, each against T};
\node[box=cGrey, anchor=north west, text width=46mm] (unk)
  at ($(meas.south west)+(0,-0.16)$)
  {\textbf{unknown}\\[1pt]
   intrinsic effects $k_{\mathrm H}/k_{\mathrm T}$,
   $k_{\mathrm D}/k_{\mathrm T}$ of the chemical step;\\
   commitment $C_{\mathrm T}=k_{\mathrm T}/k_{\mathrm{off}}$};

\node[anchor=north west, inner sep=0pt, font=\fontsize{7.6}{9}\selectfont]
  (map) at ($(unk.south west)+(0.02,-0.24)$)
  {$K_{L\mathrm T}=\dfrac{(k_L/k_{\mathrm T})\,(1+C_{\mathrm T})}
                        {1+(k_L/k_{\mathrm T})\,C_{\mathrm T}}$};
\node[ann, anchor=north west, text width=50mm]
  at ($(map.south west)+(-0.02,-0.20)$)
  {The two measurements leave one unknown commitment. Each possible commitment gives one intrinsic pair.};

% =====================================================================
% (B) the intrinsic exponent implied at every commitment
% =====================================================================
\begin{axis}[ov, name=B, width=70mm, height=50mm,
  at={(7.05cm,-0.62cm)}, anchor=north west,
  xmin=0, xmax=1, ymin=2.4, ymax=\fxTop,
  xlabel={fraction of maximum commitment, $C_{\mathrm T}(\KHT-1)$},
  ylabel={intrinsic Swain--Schaad exponent},
  xtick={0,0.25,0.5,0.75,1}, xticklabels={0,0.25,0.5,0.75,1},
  ytick={2.5,3.0,3.5,4.0,4.5,5.0},
  yticklabels={2.5,3.0,3.5,4.0,4.5,5.0},
  clip=false,
  title={(\textit{\textbf{B}})\ \ shared-tracer illustration}]
  \draw[cInk, line width=0.6pt, densely dashed]
    (axis cs:0,\fxGamma) -- (axis cs:1,\fxGamma);
  \node[ann, anchor=north east] at (axis cs:0.985,\fxGamma-0.05)
    {semiclassical reference, 3.35};

  \addplot[cPro, line width=1.1pt] table[x=m,y=e] {data/fx_exp_yadh.dat};
  \addplot[cHyd, line width=1.1pt, densely dashed] table[x=m,y=e]
    {data/fx_exp_light.dat};
  % observed exponents: the values at zero commitment
  \addplot[only marks, mark=*, mark size=1.9pt, cPro]
    table[x=m,y=e] {data/fx_obs_yadh.dat};
  \addplot[only marks, mark=*, mark size=1.9pt, cHyd]
    table[x=m,y=e] {data/fx_obs_light.dat};
  \addplot[only marks, mark=o, mark size=2.2pt, cHyd, line width=0.7pt]
    table[x=m,y=e] {data/fx_cross.dat};
  \node[ann, color=cPro, anchor=south west] at (axis cs:0.03,4.55)
    {yeast ADH, observed 3.58};
  \node[ann, color=cHyd, anchor=north west] at (axis cs:0.03,2.66)
    {ecDHFR, $5\,^{\circ}$C, observed 2.71};
  \node[ann, color=cHyd, anchor=south east] at (axis cs:0.482,3.40)
    {$C_{\mathrm T}=0.248$ (fraction $0.51$)};

  % the allowed ranges, projected onto the exponent axis
  \draw[setarr, cPro] (axis cs:1.06,\fxYadhE) -- (axis cs:1.06,\fxTop);
  \draw[cPro, fill=white, line width=0.7pt] (axis cs:1.06,\fxYadhE) circle (1.6pt);
  \draw[setarr, cHyd] (axis cs:1.13,\fxLightE) -- (axis cs:1.13,\fxTop);
  \draw[cHyd, fill=white, line width=0.7pt] (axis cs:1.13,\fxLightE) circle (1.6pt);
  \node[annc, rotate=90] at (axis cs:1.20,3.9) {range at positive commitment};
\end{axis}

\end{tikzpicture}
\end{document}
